\tocsection{\textbf{11/6/26 Lecture 30:} Orientations, Chains, and Boundaries}

Let $X$ be a topological space. We want invariants of $X$ which are homotopy invariant. We have $\pi_0$ and $\pi_1$. Using $\pi_0$, we showed $\R \not\cong \R^n$ for $n > 1$, and using $\pi_1$, we showed $\R^2 \not\cong \R^n$ for $n > 2$.\\

One way to extend to $\R^k$ is using the \emph{higher homotopy groups} $\pi_k(X, x)$ for $k \geq 0$, given by
\[ \pi_k(X, x) = \{\tt{homotopy classes of maps } S^k \to X \tt{ sending the north pole to } x\} \]

\textbf{Fact:} This has a group structure for $k \geq 1$, and for $k > 1$, these are abelian.\\

\textbf{Problem:} Hard to compute. We don't even know $\pi_k(S^2, x)$ for all $k$.\\

\textbf{Our Approach:} Simplicial Homology.\\

\textbf{Pro:} Relatively easy to compute.

\subsection*{Orientation of a Simplex}
\sectionspace

We want to distinguish
\fig{0.6}{Lecture30Fig1}{}
\textbf{First case:} Denote it with the added information (``orientation'') $(v_0, v_1, v_2)$ or a cyclic permutation of that.\\

\textbf{Second case:} Denote it with the added information (``orientation'') $(v_0, v_2, v_1)$ or a cyclic permutation of that.

\begin{definition}[\textbf{Definition 30.1:} Orientation]
    Let $A$ be a simplex other than a point, and consider orderings of its vertices. Say two orderings are \emph{equivalent} if they differ by an even permutation. The two equivalence classes are called \emph{orientations} of $A$.
\end{definition}

\begin{definition}[\textbf{Definition 30.2:} Induced Orientation]
    Say the vertices of $A$ have an orientation given by $(v_0, v_1, \ldots, v_k)$. Say $B$ is the face obtained by deleting $v_i$. We get an ordering $(v_0, \ldots, \widehat{v}_i, \ldots, v_k)$ of the vertices of $B$ (the notation means remove $v_i$).\\

    If $i$ is even, say this is the \emph{induced orientation} of $B$. If $i$ is odd, it's the opposite of the induced orientation.
\end{definition}

\begin{example}{\textbf{Example:} A Triangle}
    \fig{0.8}{Lecture30Fig2}{}
    The orientation of $A$ is $(v_0, v_1, v_2)$. Then, $(v_1, v_2)$ is right, $(v_0, v_2)$ is wrong so $(v_2, v_0)$ is right, and $(v_0, v_1)$ is right:
    \fig{0.8}{Lecture30Fig3}{}
\end{example}

Say $(G, +)$ is abelian.

\begin{definition}[\textbf{Definition 30.3:} Generating Sets]
    A subset $X$ of $G$ is a \emph{generating set} if every $g \neq 0$ can be written as
    \[ g = n_1x_1 + \cdots + n_kx_k \]
    where $n_i \in \Z \setminus \{0\}$, $x_i \in X$, and $x_i \neq x_j$ if $i \neq j$.
\end{definition}

\begin{definition}[\textbf{Definition 30.4:} Free Abelian Groups]
    A generating set is \emph{free} if such a representation as above is unique. We then call $G$ a \emph{free abelian group}.
\end{definition}

\begin{definition}[\textbf{Definition 30.5:} Abelian Words]
    Given a set $X$, an \emph{abelian word} is a formal finite sum $\sum n_ix_i$, $n_i \in \Z$, $x_i \in X$. It's \emph{reduced} if $n_i \neq 0$ and $x_i \neq x_j$ if $i \neq j$. The reduced abelian words form a group $FA(X)$ generated by $X$.
\end{definition}

Now, if $G$ is an abelian group and $X$ is a set of generators, we get a homomorphism
\[ \phi : FA(X) \to G \quad \tt{by} \quad \phi\left(\sum n_ix_i\right) = \sum n_ix_i \]
$\phi$ is onto, so let $N = \ker(\phi)$. Then, we see $G \cong FA(X)/N$, so we can talk about presentations and finite presentations.\\

Given an abelian group $G$, set
\[ T = \{g \in G \col ng = 0 \tt{ for some } n \neq 0\} \]
called the \emph{torsion subgroup}. Note that then, $G/T$ is torsion free (it has a trivial torsion subgroup).\\

If $G$ is finitely generated, so is $G/T$, and it turns out we can write $G/T \cong \Z^n$ for some $n$. We call $n$ the \emph{rank}, $\tt{rk}(G)$, of $G$.\\

Note that we can write $G \cong T \oplus G/T$, but the isomorphism is not unique.

\subsection*{Simplicial Complexes}
\sectionspace

Let $K$ be a simplicial complex.

\begin{definition}[\textbf{Definition 30.6:} Chain Groups]
    Let $C_q(K)$ be the free abelian group generated by the oriented $q$-simplices of $K$, modulo the objects of the form $\sigma + \tau$ if $\sigma$ and $\tau$ are the same simplex with opposite orientations. If $q = 1$, then as there is only one orientation of a vertex, $C_{0}(K)$ is just the free abelian group on the vertices.
\end{definition}

\begin{definition}[\textbf{Definition 30.7:} Chains]
    An element of $C_q(K)$ is called a \emph{$q$-chain}.
\end{definition}

A $q$-chain is something of the form $\sum \lambda_i\sigma_i$, where $\lambda_i \in \Z$ and $\sigma_i$ is an oriented $q$-simplex in $K$. Remember $(-\lambda)\sigma = \lambda(-\sigma)$, where $\lambda \in \Z$ and $-\sigma$ denotes $\sigma$ with the opposite orientation.\\

It is clear that $C_q(K) \cong \Z^{n_q}$, where $n_q$ is the number of $q$-simplices in $K$.

\begin{definition}[\textbf{Definition 30.8:} Boundary]
    If the oriented $q$-simplex has vertices $(v_0, \ldots, v_q)$, then the \emph{boundary} is
    \[ \partial(v_0, \ldots, v_q) = \sum_{i = 0}^{q} (-1)^{i}(v_0, \ldots, \widehat{v}_i, \ldots, v_q) \in C_{q - 1}(K) \]
\end{definition}

\begin{example}{\textbf{Example:} Boundary of a Triangle}
    \fig{0.8}{Lecture30Fig4}{}
    \begin{align*}
        \partial(v_0, v_1, v_2) &= (v_1, v_2) - (v_0, v_2) + (v_0, v_1)\\
        &= (v_1, v_2) + (v_2, v_0) + (v_0, v_1) \in C_1(K)
    \end{align*}
    \fig{0.8}{Lecture30Fig5}{}
\end{example}

Now, we can extend to a map $\partial : C_q(K) \to C_{q - 1}(K)$ by
\[ \partial\left(\sum n_i\sigma_i\right) = \sum n_i(\partial\sigma_i) \]
If $q = 0$, we define $C_{-1}(K) = \{0\}$, so for $q < 1$, $\partial$ is the zero map.

\begin{definition}[\textbf{Definition 30.9:} Cycles and Boundaries]
    Let
    \[ Z_q(K) = \ker(\partial : C_q(K) \to C_{q - 1}(K)) \]
    be the set of \emph{$q$-cycles}, and let
    \[ B_q(K) = \operatorname{Im}(\partial : C_{q + 1}(K) \to C_q(K)) \]
    be the set of \emph{$q$-boundaries}.
\end{definition}

\begin{example}{\textbf{Example:} Edge Paths}
    Say $(v_0, \ldots, v_{\ell})$ is an edge path (with $v_i \neq v_{i + 1}$). Let
    \[ \alpha = (v_0, v_1) + (v_1, v_2) + \cdots + (v_{\ell - 1}, v_{\ell}) \in C_1(K) \]
    We see
    \begin{align*}
        \partial\alpha &= ((v_1) - (v_0)) + ((v_2) - (v_1)) + \cdots + ((v_{\ell}) - (v_{\ell - 1}))\\
        &= (v_{\ell}) - (v_0)
    \end{align*}
    \fig{0.6}{Lecture30Fig6}{}
    If our edge path is a loop, with $v_{\ell} = v_0$, then $\partial\alpha = 0$, so $\alpha$ is a $1$-cycle, so $\alpha \in Z_1(K)$.
\end{example}
